\documentclass[12pt]{article}
\usepackage[margin=1in]{geometry}
\usepackage{amsmath,booktabs,threeparttable}

\begin{document}

\providecommand{\StructuralRobustnessPath}{appendix/structural-robustness}

\subsection{Robustness to Heavy-Tailed Intrinsic Motivation}
\label{app:student-t-types}

The baseline specification assumes that intrinsic motivation $v_i$ is
standard normal. We re-estimate the model assuming instead that $v_i$ follows
a centered Student-$t$ distribution with five degrees of freedom, scaled to
have unit variance. The idiosyncratic decision shock $u_i$ remains Gaussian.
Consequently, this exercise changes the tail thickness of intrinsic
motivation while preserving its location and variance and leaving the rest of
the structural model unchanged.

The convolution of a Student-$t$ type with Gaussian decision noise has no
convenient closed-form distribution. We evaluate it using the exact
normal-variance-mixture representation of the Student-$t$ and a deterministic
12-node generalized Gauss--Laguerre quadrature rule. The nodes and weights are
fixed numerical integration constants, not estimated latent classes. Across
the range used by the model, the approximation error is below $3\times10^{-4}$
for the Student-$t(5)$ distribution function and below $2\times10^{-4}$ for
its density. As an implementation check, selecting the Gaussian branch
recovers the frozen baseline likelihood and gradient exactly at the audit
points.

The Student-$t$ fit used four chains with 400 warmup and 400 retained draws per
chain. The maximum split $\widehat R$ was 1.008, minimum bulk and tail effective
sample sizes were 475 and 549, respectively, and there were no divergences or
maximum-treedepth transitions.

Table~\ref{tab:student-t-multipliers} compares the resulting social
multipliers. Heavy-tailed types modestly attenuate the estimated multipliers,
with reductions of approximately 1--10 percent that are generally larger at
longer distances. The qualitative pattern is unchanged. The multiplier
remains above one for Control and Calendar and below one for Ink and Bracelet;
for Calendar, the Student-$t$ credible intervals include one at 1.5 and 2.5
km. Thus the conclusion that social image amplifies the distance response in
the less visible arms and mitigates it in the signal arms does not depend on a
Gaussian latent-type distribution.

\begin{table}[htbp]
  \centering
  \small
  \caption{Social multipliers under Gaussian and heavy-tailed intrinsic motivation}
  \label{tab:student-t-multipliers}
  \begin{threeparttable}
    \begin{tabular}{lccc}
\toprule
Treatment & 500m & 1500m & 2500m \\
\midrule
\multicolumn{4}{l}{\textit{Panel A: Gaussian intrinsic motivation}} \\
Bracelet & 0.89 & 0.83 & 0.78 \\
 & (0.70, 1.04) & (0.69, 0.95) & (0.67, 0.90) \\
Calendar & 1.26 & 1.22 & 1.21 \\
 & (1.11, 1.47) & (1.05, 1.46) & (1.01, 1.52) \\
Ink & 0.92 & 0.86 & 0.81 \\
 & (0.80, 1.03) & (0.74, 0.97) & (0.70, 0.94) \\
Control & 1.43 & 1.48 & 1.57 \\
 & (1.24, 1.73) & (1.24, 1.87) & (1.26, 2.03) \\
\addlinespace
\multicolumn{4}{l}{\textit{Panel B: Student-$t(5)$ intrinsic motivation}} \\
Bracelet & 0.88 & 0.78 & 0.71 \\
 & (0.70, 1.03) & (0.65, 0.90) & (0.60, 0.82) \\
Calendar & 1.23 & 1.13 & 1.11 \\
 & (1.09, 1.42) & (0.98, 1.34) & (0.91, 1.37) \\
Ink & 0.87 & 0.79 & 0.73 \\
 & (0.75, 0.97) & (0.67, 0.90) & (0.60, 0.86) \\
Control & 1.38 & 1.37 & 1.45 \\
 & (1.20, 1.66) & (1.17, 1.70) & (1.18, 1.87) \\
\bottomrule
\end{tabular}

    \begin{tablenotes}\footnotesize
      \item \textit{Notes:} Entries are posterior medians of the social
      multiplier evaluated at the indicated distance. Parentheses contain 95
      percent posterior credible intervals. The Student-$t$ specification has
      five degrees of freedom and is scaled to unit variance; its convolution
      with Gaussian decision noise is evaluated using deterministic 12-node
      quadrature.
    \end{tablenotes}
  \end{threeparttable}
\end{table}

\end{document}
